\documentclass[11pt,a4paper,roman]{moderncv}
\moderncvstyle{banking}
\moderncvcolor{purple}
\nopagenumbers{}
\usepackage[utf8]{inputenc}
\usepackage[T1]{fontenc}
\usepackage{lmodern}
\usepackage[scale=0.96,top=0.5cm,bottom=0.5cm,left=0.75cm,right=0.75cm]{geometry}
\renewcommand{\baselinestretch}{0.93}
\setlength{\hintscolumnwidth}{2.25cm}
\setlength{\footskip}{18pt}

\name{Wahid}{SAMER}
\address{Cergy (95), France}{}{}
\mobile{07 53 10 95 74}
\email{samer.ahmed.wahid@gmail.com}
\social[linkedin]{Wahid SAMER}
\social[github]{wahidrt}
\title{Alternance M2 - Actuariat / Data \& Modélisation des risques}

\begin{document}
\makecvtitle
\vspace{-3.3em}
\begin{center}
\parbox{0.97\textwidth}{\centering\small\itshape
Étudiant en mathématiques financières, orienté \textbf{actuariat, analyse de données et modélisation des risques}, avec une solide base en probabilités, statistiques, simulation et programmation. Recherche d'une \textbf{alternance M2 de 12 mois à partir du 21 septembre 2026}. Mobilité : \textbf{France, Luxembourg, Suisse, Belgique}.}
\end{center}
\vspace{0.15em}

\section{Formation}
\cventry{2025 -- 2027}{Master Mathématiques Appliquées à l'Ingénierie Financière}{CY Cergy Paris Université / CY Tech}{Cergy}{}{\textbf{M1 :} probabilités, processus stochastiques, simulation stochastique, statistiques avancées, optimisation, méthodes numériques et gestion des risques.\\
\textbf{M2 2026--2027 :} grands risques et valeurs extrêmes, séries temporelles et machine learning avec Python, model calibration and simulation, portfolio management et fixed income.}
\cventry{2021 -- 2025}{Licence Mathématiques et Applications}{Université d'Évry Paris-Saclay}{Évry}{}{Probabilités, statistique, mesure et intégration, analyse, algèbre, analyse numérique et programmation.}

\section{Projets ciblés}
\cventry{2026}{Python}{Modélisation et suivi des risques d'un portefeuille}{Projet M1 - équipe de 4}{}{Calcul du P\&L et des expositions, suivi des coûts et du drawdown, indicateurs Sharpe/Sortino, stress de paramètres et analyse de scénarios sur données financières.}
\cventry{2026}{R}{Analyse statistique de données socio-économiques}{Projet universitaire}{}{Préparation et croisement de données sur l'espérance de vie et le PIB ; analyse descriptive et visualisation avec \texttt{dplyr}, \texttt{tidyr} et \texttt{ggplot2}.}
\cventry{2026}{Python / \LaTeX}{Black--Scholes, Heston et sensibilités}{Mémoire M1}{}{Modélisation stochastique, simulation, méthodes numériques et analyse des sensibilités ; comparaison d'un modèle à volatilité constante avec une volatilité stochastique.}
\cventry{2025}{Excel / VBA}{Outil de calcul et scénarios de risque}{Projet M1}{}{Automatisation du pricing de Call/Put, calcul des sensibilités $\Delta,\Gamma,\Theta,\nu,\rho$ et scénarios de stress sur la volatilité.}

\section{Compétences}
\cvitem{Actuariat / risque}{Probabilités, statistique, processus stochastiques, simulation, valeurs extrêmes (M2), analyse de scénarios, sensibilités et risque de portefeuille.}
\cvitem{Data}{Python (Pandas, NumPy, SciPy, Scikit-learn, Matplotlib), R (dplyr, tidyr, ggplot2), SQL.}
\cvitem{Outils}{Excel / VBA, Git / GitHub, \LaTeX, Power BI ; restitution et visualisation de résultats.}
\cvitem{Langues}{Français : courant ; Arabe : langue maternelle ; Anglais : intermédiaire.}

\section{Expérience \& Engagement}
\cventry{Mai -- juin 2025}{Professeur de mathématiques - stagiaire}{Lycée Pierre Mendès France}{Vitrolles}{}{Préparation et transmission de notions de niveau Terminale ; rigueur, pédagogie et communication scientifique.}
\cventry{2019}{Technicien aéronautique - stagiaire}{G1 Aviation}{Gap}{}{Assemblage et maintenance dans un environnement exigeant en précision, procédures et sécurité.}
\cvitem{Associatif}{Bénévole aux \textbf{Restaurants du Cœur} pendant près de 3 ans : solidarité, accueil et travail d'équipe.}

\end{document}
